\section{A divisor itinerary and its small observations}\label{sec:itinerary}

Fix $0<\eta<1/4$, put $\xi=\eta^3$, and fix a threshold
$P_0\ge2$, independent of the height $T$.  We count nonendpoint
points $x=a/q\in C$ in lowest terms, with $(q,3)=1$ and
$T/2<q\le T$.  We can discard $q=1$.  Tripling modulo one permutes
the points of each such denominator: it shifts their periodic ternary
expansions, and its $n$th iterate is the identity, where
$n=\operatorname{ord}_q(3)$.  We use the following published consequence
of Bourgain's theorem: for every $\delta>0$, $n\ll_\delta q^\delta$
for represented $q$ coprime to $3$ \cite[\S6.2]{RSTW2020}. Thus
$n=T^{o(1)}$ uniformly here.  No bound of logarithmic order is needed.

For $h\mid n$ let
\[
 b(h)=(q,3^h-1),\qquad d(h)=q/b(h),\qquad
 W_h=\{3^h x\}-x.
\]
The last variable is a \emph{real} number in $(-1,1)$, whose class in
$\mathbb R/\mathbb Z$ has exact denominator $d(h)$ (the denominator
of zero is one). Initially $h=n$, $b(n)=q$ and $W_n=0$.
Write $m_{>P_0}$ for the factor of $m$ that remains on removing all
prime powers based at primes $\le P_0$.

\begin{lemma}[The denominator itinerary]\label{lem:itinerary}
There is a deterministic rule depending only on $q,T,\eta,P_0$ for
choosing a chain $n=h_0>\cdots>h_J=h_*$ of divisors, with
$h_*$ supported on primes $\le P_0$. The increments
$u_i=b(h_{i-1})/b(h_i)$ and terminal factor $b=b(h_*)$ satisfy
$b\prod_i u_i=q$. Each retained step is of one of two types:
\begin{enumerate}
\item A large step $h_{i-1}=ph_i$, where $p>P_0$ is prime and
$u_i>T^\eta$. Put $L=(h_{i-1})_{>P_0}$. Of the prime powers
$r^e\parallel b(h_{i-1})$ with
$(\operatorname{ord}_{r^e}3)_{>P_0}<L^\xi$, those not surviving
in full in $b(h_i)$ have product at most $T^\xi$.
\item A group of consecutive small steps, of combined increment at
most $T^{2\eta}$; its terminal and initial lags are both retained.
\end{enumerate}
There are $J\le 3/\eta+1$ steps.  The number of possible labels
consisting of $n$, the retained lags and types, and dyadic size
bins for all $u_i$ and $b$, is $T^{o(1)}$ with parameters fixed.
\end{lemma}
\begin{proof}
At lag $g$ with $L=g_{>P_0}>1$, choose a candidate $p\mid L$ at random
with probabilities $v_p(L)\log p/\log L$. A full component
$r^e\parallel b(g)$ of order $d=\operatorname{ord}_{r^e}3$ is not
entirely in $b(g/p)$ exactly if $d\nmid g/p$. As $d\mid g$, this
requires $v_p(d)=v_p(g)$. Its probability is at most
$\log(d_{>P_0})/\log L$. By linearity the expected sum of
$\log(r^e)$ over lost components with $d_{>P_0}<L^\xi$ is at most
$\xi\log b(g)\le\xi\log T$, even charging partially lost components
in full. Take a minimizer, breaking ties by the smallest $p$, and
repeat until $L=1$.

Keep each raw increment $>T^\eta$ unchanged. In each intervening
run, group consecutive increments until their product first reaches
$T^\eta$, then start again; keep a nonempty trailing group even if
its product is one. A completed group has product at most $T^{2\eta}$.
All increments and the terminal multiply to $q$, by cancellation.
There are at most $1/\eta$ large increments, at most $1/\eta$
completed groups, and at most $1+1/\eta$ trailing groups.
The retained lags are $\le n=T^{o(1)}$, and there are $O(\log(2T))$
bins per slot. All labels are functions of the denominator.
\end{proof}

\begin{lemma}[Fibres at fixed past]\label{lem:trace-fibres}
For a retained step $g\to h$, set $B=(3^g-1)/(3^h-1)$.
At fixed real $W_g$ and exact $u=b(g)/b(h)$, the torus values of
$W_h$ lie in a single translate of $G_u=(\tfrac1u\mathbb Z)/\mathbb Z$,
and $u\mid B$. Each such value has at most two real lifts.
At the terminal, fixing $W_{h_*}$ and $b$ puts $x$ in one translate
of $G_b$.
\end{lemma}
\begin{proof}
Modulo one, $B W_h=W_g$, and $d(h)/(B,d(h))=d(g)$. Thus
$(B,d(h))=u$; $d(g)$ is known from $W_g$, so $d(h)=ud(g)$ is
common to all the candidates. Multiplication by $B$ on $G_{d(h)}$
has kernel $G_u$. At the terminal $q=b\,d(h_*)$ is fixed and the
kernel of $3^{h_*}-1$ on $G_q$ is $G_b$.
\end{proof}

\begin{lemma}[A joint entropy bound for the small observations]\label{lem:small-traces}
Fix a nonempty label of Lemma~\ref{lem:itinerary}, and let $X$ be
uniform on all its points. Write $W_i=W_{h_i}(X)$.
There is finite side information $E$ such that, in natural logarithms,
\[
 H(E)+\sum_{i\ {\rm small}} H(W_i\mid E,W_0,\ldots,W_{i-1})
 \ \le\ O\big(\eta\log(2/\eta)\big)\log T+O_\eta(1),
\]
with constants independent of $P_0$. No conditional distribution on a
real past is required to be invariant under tripling.
\end{lemma}
\begin{proof}
Write $X=\sum_{t\ge0}2w_t3^{-t-1}$ with $w_t\in\{0,1\}$ indexed
periodically modulo $n$. The ternary trits of
$Z_i=(1+W_i)/2$ are precisely $1+w_{t+h_i}-w_t$. Indeed their
series sums to $Z_i$ without any carry. These trits are unique:
$Z_i\in(0,1)$ has denominator prime to $3$, hence is not a triadic
endpoint.

For a small step let $U_i$ be the public upper bin endpoint, taken
$\le2T^{2\eta}$. Given $W_{i-1}$, the first
$M_i=\lceil 2\log_3 U_i\rceil+3$ trits specify $W_i$,
\emph{without} being given $u_i$. In fact the torus difference of
two candidates is killed by $B$ in Lemma~\ref{lem:trace-fibres}
and lies in $G_{\operatorname{lcm}(d,d')}$, where
$d,d'$ are their child denominators. Its order divides
\[
 (B,\operatorname{lcm}(d,d'))
 =\operatorname{lcm}((B,d),(B,d'))
 =\operatorname{lcm}(u,u')\le U_i^2.
\]
Distinct real $Z_i$ are thus separated by at least $1/(2U_i^2)$,
including when the torus values agree.

Let $M$ be the maximal such $M_i$ (zero if there are no small
steps). Say that $F_i(t)$ holds if the bits in the entire class
$t+h_i\mathbb Z$ modulo $n$ are constant.  The events decrease as
$i$ increases, and $F_0(t)$ always holds. For $0\le t<M$ record in
$E$ the bit $w_t$ and the first index $i$ at which $F_i(t)$ fails,
or a marker if it never fails. The indices include the large steps.
Then $H(E)\le M\log(2(J+1))$.

From $W_{i-1}$ one knows the differences
$w_{s+h_{i-1}}-w_s$ for every $s$. On each nonconstant
$h_{i-1}$-cycle they recover all the bits: a nonzero difference
fixes both endpoints, after which one integrates around the cycle.
At a site $t<M_i$ of a small observation, its trit is known if
$F_i(t)$ holds. It is also known if $F_{i-1}(t+h_i)$ fails, because
then the target bit is recovered and the source bit is prepaid.
The remaining event
$F_{i-1}(t+h_i)\setminus F_i(t)$ is known given $E,W_{i-1}$:
the first event is read from the old difference cycle and the
second from the recorded index. On that event only the target bit
is unknown. The original uniform distribution is shift-invariant,
since its label depends only on $q$. With $\theta_i=\Pr(F_i(0))$,
the remaining event has probability $\theta_{i-1}-\theta_i$:
indeed $F_i(t)=F_i(t+h_i)\subseteq F_{i-1}(t+h_i)$.
Subadditivity of conditional entropy and the trit reconstruction
give the upper bound $M_i(\theta_{i-1}-\theta_i)\log 2$ for this
step. Across all small steps the sum is at most $M\log 2$, since
the omitted drops at large steps are nonnegative. Now
$M\le4\eta\log_3 T+O(1)$ and $J=O(1/\eta)$.
\end{proof}

\paragraph{The arithmetic available after conditioning.}
At a large raw step $g=ph\to h$ put $L=g_{>P_0}$.
Every increment $u$ splits into (i) a product of size $\le T^\xi$
from the low components of Lemma~\ref{lem:itinerary}, (ii) components
based at primes dividing $n$ not already counted, and (iii) whole
powers $r^e\parallel q$, $r\nmid 3n$, with
$d_r=\operatorname{ord}_r3\mid g$ and $(d_r)_{>P_0}\ge L^\xi$.
Indeed the kernel of reduction of units modulo $r^e$ to units
modulo $r$ has $r$-power order, whereas
$\operatorname{ord}_{r^e}3\mid n$ for $r^e\mid q$. For $r\nmid n$
this gives equality of the orders at every depth up to $e$; the
whole component survives or disappears. At the terminal the analogous
description has no low factor and $d_r\mid h_*$ is supported on
primes $\le P_0$. For fixed $n$, any products $\le T$ supported on
primes dividing $n$ have only $T^{o(1)}$ names: for each fixed
$v>0$ their number is at most
$T^v\prod_{r\mid n}(1-r^{-v})^{-1}
 \le T^v\exp(O_v(\log(2n)))$; then take $v$ arbitrarily small.

\begin{lemma}[Names of subproducts]\label{lem:itinerary-names}
Fix the large-step or terminal order restrictions in the preceding
paragraph. Let $\mathcal R$ be the set of allowed primes $r\nmid3n$.
For $1/2<s<1$, products of distinct bases in $\mathcal R$ of size
$\le R$ have $O_s(R^s)$ names. For $a>0$ restrict further to
$d_r<r^a$; their base products have $O_{a,v}(R^v)$ names for
every $v>a$. Finally products $\le R$ of powers $r^e$ with $e>A$
have $O_{A,v}(R^v)$ names whenever $(A+1)v>1/2$, for an integer
$A\ge1$. Constants here may also depend on $P_0,\xi$, but not on
the lags or $T$. These bounds hold for arbitrary subproducts, even
when selected using the actual point sets at a fixed past.
\end{lemma}
\begin{proof}
At exact order $d$ there are $O(d)$ primes, since their product
divides $3^d-1$. They lie among $1+jd$ for distinct integers $j\ge1$.
Thus their total reciprocal $s$-mass is
$O_s(d^{-s}\sum_{j\le O(d)} j^{-s})=O_s(d^{1-2s})$,
by taking the smallest possible indices. Under $d<r^a$ their
reciprocal $v$-mass is $O(d^{1-v/a})$. Either estimate is
$O(d^{-c})$ for some fixed $c>0$.
Writing $d=ke$ with $k$ supported on primes $\le P_0$ and
$e\mid L$, $e\ge L^\xi$, the sum over the permissible orders of
$d^{-c}$ at a large step is at most
\[
 \prod_{\ell\le P_0\ {\rm prime}}(1-\ell^{-c})^{-1}
 \,\tau(L)L^{-\xi c}=O_{P_0,\xi,c}(1).
\]
At the terminal only the finite Euler product is needed. The
divisor bound makes the estimate uniform even as $L$ varies.
Consequently the Euler product of $1+r^{-s}$, or $1+r^{-v}$ for
the restricted low-order bases, is bounded. For large exponents,
$\sum_{e>A}r^{-ev}\le r^{-(A+1)v}/(1-2^{-v})$; the bounded
base-mass estimate applies with any
$1/2<s<\min(1,(A+1)v)$. Multiplying the corresponding local
factors and applying $1_{m\le R}\le(R/m)^v$ (or exponent $s$)
proves the assertions. None of these bounds enumerates new lag
paths for the subproducts.
\end{proof}

In particular, for a bounded multiplicity $e$ one may group
remaining unramified components by $e$ and bin their \emph{radical}
products: knowing the radical in that slot determines its $e$th
power. Bins for such subproducts at a fixed past are for counting
only. Adding them to the original label before the entropy
calculation would require a new proof of rotation invariance.
